\section{Introducción: ¿por qué se necesita la compresión de contexto?}

Cuando se usa un Coding Agent como Codex durante sesiones de desarrollo largas, el contexto de la conversación se va acumulando continuamente. Cuando la longitud del contexto se acerca al límite de la Context Window del modelo, si no se procesa, el Agent ya no podrá recibir información nueva. La \textbf{compresión de contexto} (Compact) es precisamente el mecanismo diseñado para resolver este problema: bajo la premisa de conservar la información clave, comprime el historial extenso de la conversación en un resumen conciso, de modo que el Agent pueda "continuar sin interrupciones" con el trabajo posterior.

Esta entrega se basa principalmente en un análisis inverso, publicado en un conocido hilo de Twitter, del mecanismo de Compact de Codex, combinado con la lectura del código fuente de Codex.

\begin{knowledgebox}{Pregunta central}
¿Cómo funciona la compresión de contexto (Compact) de Codex? ¿Cuál es el System Prompt detrás de ella? Después de la compresión, ¿cómo se hace el traspaso (Handoff) al siguiente LM?
\end{knowledgebox}

